% Source-derived chapter 01: Introduction
% Original source: geometric-local-systems-general-curves-isomonodromy
\chapter{Introduction}
\label{geometric-local-systems-general-curves-isomonodromy-chapter-01}
\source{geometric-local-systems-general-curves-isomonodromy}

\begin{remark}[Source section]
\label{geometric-local-systems-general-curves-isomonodromy-chapter-01-source}
\source{geometric-local-systems-general-curves-isomonodromy}
\source{geometric-local-systems-general-curves-isomonodromy}
This chapter reproduces the corresponding section of the retrieved
LaTeX source, including its definitions, statements, examples, and
proofs. It has not yet been translated into Lean.
\notready
\end{remark}

% The source section title was: Introduction
\label{section:introduction}
\source{geometric-local-systems-general-curves-isomonodromy}
\subsection{Overview}
\label{subsection:overview}
\source{geometric-local-systems-general-curves-isomonodromy}
We work over the complex numbers $\mathbb{C}$. 
The main result of this paper, \autoref{theorem:very-general-VHS}, is that an
analytically very general $n$-pointed curve of genus $g$ 
(defined in \autoref{definition:general})
does not carry any
non-isotrivial polarizable integral variations of Hodge structure of rank less than $2\sqrt{g+1}$. In particular, an
analytically very general $n$-pointed curve of genus $g$ carries no geometric local
systems of rank less than $2\sqrt{g+1}$ with infinite monodromy, as we show in
\autoref{corollary:geometric-local-systems}. This is a strong restriction on the
topology of smooth proper maps to an analytically very general curve, and, as pointed out by H\'el\`ene Esnault, contradicts conjectures of Esnault-Kerz \cite[Conjecture 1.1]{esnault2021local} and Budur-Wang \cite[Conjecture 10.3.1]{budur2020absolute}, as explained in \autoref{corollary:non-density-of-geometric-local-systems}. 

The above results rely on an analysis of stability properties of isomonodromic deformations of flat vector bundles with regular singularities, and require correcting a number of errors in the literature on this topic.
We next state our main results on stability properties of isomonodromic
deformations of flat vector bundles.
Let $C_0$ be the central fiber of a family of curves $\mathscr{C}\to \Delta$
with $\Delta$ a contractible domain, and let $(E_0, \nabla_0)$ be a vector
bundle with flat connection
on $C_0$. Recall that, loosely speaking, the isomonodromic deformation of $(E_0, \nabla_0)$  is the deformation $(\mathscr{E}, \nabla)$  of $(E_0, \nabla_0)$ to $\mathscr{C}/\Delta$, such that the monodromy of the connection is constant.

In \autoref{corollary:counterexample}, we construct a flat vector bundle on a smooth proper curve over $\mathbb{C}$, whose isomonodromic deformations 
to a nearby curve
are never semistable. 
(See \autoref{definition:nearby} for precise definitions.)
The construction arises from the ``Kodaira-Parshin trick," and contradicts  earlier claimed theorems of Biswas, Heu, and Hurtubise (\cite[Theorem 1.3]{BHH:logarithmic},
\cite[Theorem 1.3]{BHH:irregular}, and
\cite[Theorem 1.2]{BHH:parabolic}), which imply that such a construction is impossible. See \autoref{remark:bhh-error} for a discussion of the errors in those papers. 

As a complement to this example, we show in
\autoref{theorem:hn-constraints-parabolic} that any logarithmic flat vector
bundle admits an isomonodromic deformation to a nearby curve which is \emph{close} to semistable,
in a suitable sense,
and moreover is (parabolically) semistable if the rank is small compared to the genus of the curve.
While our results contradict those of \cite{BHH:logarithmic, BHH:irregular,
BHH:parabolic}, our methods owe those papers a substantial debt. 
Biswas, Heu,
and Hurtubise 
pitch the question of isomonodromically deforming a vector bundle to a
semistable vector bundle (see \autoref{question:bhh})
as an
analogue of Hilbert's 21st problem, also known as the Riemann-Hilbert problem.

The semistability property of \autoref{theorem:hn-constraints-parabolic} is also the main
input to our Hodge-theoretic main results, mentioned above.
The applications to polarizable variations of Hodge structures come from the fact that flat vector bundles underlying polarizable variations are rarely (parabolically) semistable, due to well-known curvature properties of Hodge bundles.



\subsection{Main Hodge-theoretic results}
\label{subsection:hodge-intro}
\source{geometric-local-systems-general-curves-isomonodromy}
Results from this subsection, \autoref{subsection:hodge-intro}, as well as the
next, \autoref{subsection:isomonodromy-intro}, will be proven later in the
paper, as detailed in \autoref{subsection:organization}.
For convenience, throughout the paper, out main results will primarily be stated for hyperbolic curves.
\begin{definition}
\source{geometric-local-systems-general-curves-isomonodromy}
\notready
	\label{definition:hyperbolic}
\source{geometric-local-systems-general-curves-isomonodromy}
	Let $C$ be a curve over $\mathbb C$ of genus $g$ and $D \subset C$ a
	reduced effective divisor of degree $n$. 
	Call $(C, D)$ {\em hyperbolic} if $C$ is a smooth proper connected curve
	and
	either $g \geq 2$ and $n\geq 0$, $g = 1$ and $n > 0$, or $g =0$ and $n > 2$.
	We call an $n$-pointed curve $(C, x_1, \ldots, x_n)$ {\em hyperbolic} if
	$(C, x_1 +\cdots + x_n)$ is hyperbolic.
\end{definition}
\begin{remark}
\source{geometric-local-systems-general-curves-isomonodromy}
\notready
	\label{remark:}
\source{geometric-local-systems-general-curves-isomonodromy}
	Equivalently, $(C,D)$ is hyperbolic if and only if it has no infinitesimal automorphisms,
i.e., $H^0(C, T_C(-D)) = 0$.
\end{remark}

We will also work with the following analytic notion of a (very) general general point.

\begin{definition}
\source{geometric-local-systems-general-curves-isomonodromy}
\notready
	\label{definition:general}
\source{geometric-local-systems-general-curves-isomonodromy}
	A property holds for an {\em analytically general} point of a complex orbifold $X$, 
	if there exists a nowhere dense
	closed analytic subset $S \subset X$ so that the property holds
	on $X - S$.
	We say that a property holds for an
	{\em analytically very general} point if, locally on $X$, there exists a countable collection of nowhere dense closed analytic subsets such that the property holds on the complement of their union. If $\mathscr{M}_{g,n}$ is the analytic moduli stack of
$n$-pointed curves of genus $g$, we say that a property holds for an
analytically (very) general $n$-pointed curve if it holds for an analytically
(very) general point of $\mathscr{M}_{g,n}$.
\end{definition}
\begin{remark}
\source{geometric-local-systems-general-curves-isomonodromy}
\notready\label{remark:}
\source{geometric-local-systems-general-curves-isomonodromy}
	From the definition, it may appear that ``analytically very general'' is
	a local notion, while ``analytically general'' is a global notion.
	However, being ``analytically general'' also has the following
	equivalent local definition, which is more similar to the definition of
	``analytically very general'':
	locally on $X$, there exists a nowhere dense closed analytic subset such
	that the property holds on the complement of this subset.
\end{remark}

The main geometric consequence of this work is the following constraint on the
rank of non-isotrivial polarizable variations of Hodge structure (defined in
\autoref{section:hodge-theoretic-preliminaries}) on an analytically very general curve:
\begin{theorem}
\source{geometric-local-systems-general-curves-isomonodromy}
\notready\label{theorem:very-general-VHS}
\source{geometric-local-systems-general-curves-isomonodromy}
	Let $K$ be a number field with ring of integers
$\mathscr{O}_K$. Suppose 
$(C, x_1, \cdots, x_n)$ 
is an analytically very general $n$-pointed
hyperbolic curve of genus $g$, and $\mathbb{V}$ is a $\mathscr{O}_K$-local system on $C\setminus\{x_1, \cdots, x_n\}$ with infinite monodromy. 
Suppose additionally that for each embedding $\iota: \mathscr{O}_K\to \mathbb{C}$, $\mathbb{V}\otimes_{\mathscr{O}_K, \iota}\mathbb{C}$ underlies a polarizable complex variation of Hodge structure. 
Then, $$\on{rk}_{\mathscr{O}_K}(\mathbb{V})\geq 2\sqrt{g+1}.$$
\end{theorem}
\begin{remark}
\source{geometric-local-systems-general-curves-isomonodromy}
\notready
	\label{remark:}
\source{geometric-local-systems-general-curves-isomonodromy}
	Note that a result analogous to \autoref{theorem:very-general-VHS} does
	not hold for variations without an underlying $\mathscr{O}_K$-structure. 
	Indeed, every smooth proper curve of genus at least $2$ admits a polarizable complex  variation of Hodge structure of rank $2$ with infinite monodromy, arising from uniformization (see e.g.~\cite[bottom of p.~870]{simpson:constructing-VHS}).
\end{remark}


Let $X$ be a smooth variety. We say a complex local system $\mathbb{V}$ on $X$ is \emph{of geometric origin} if there exists a dense open $U\subset X$, and a smooth proper morphism $f: Y\to U$ such that $\mathbb{V}|_U$ is a direct summand of $R^if_*\mathbb{C}$ for some $i\geq 0$. As local systems of geometric origin satisfy the hypotheses of \autoref{theorem:very-general-VHS}, we have:
\begin{corollary}
\source{geometric-local-systems-general-curves-isomonodromy}
\notready\label{corollary:geometric-local-systems}
\source{geometric-local-systems-general-curves-isomonodromy}
Let $(C, x_1, \cdots, x_n)$ be an analytically very general hyperbolic $n$-pointed curve of genus $g$. If $\mathbb{V}$ is a local system on $C\setminus\{x_1, \cdots, x_n\}$ of geometric origin and with infinite monodromy, then $\dim_{\mathbb{C}}\mathbb{V}\geq 2\sqrt{g+1}$.
\end{corollary}
We will prove \autoref{corollary:geometric-local-systems} 
in \autoref{subsubsection:geometric-origin-proof}.
As a consequence of \autoref{corollary:geometric-local-systems} 
we obtain the following concrete geometric corollary:
\begin{corollary}
\source{geometric-local-systems-general-curves-isomonodromy}
\notready
	\label{corollary:abelian-schemes}
\source{geometric-local-systems-general-curves-isomonodromy}
	If $(C, x_1, \ldots, x_n)$ is an analytically
very general hyperbolic $n$-pointed genus $g$ curve, then any non-isotrivial
abelian scheme over $C\setminus\{x_1, \cdots, x_n\}$  has relative dimension at least $\sqrt{g+1}$.
	Similarly, any relative smooth proper curve over $C\setminus\{x_1, \cdots, x_n\}$
	has genus at least $\sqrt{g+1}$.
\end{corollary}
We will prove \autoref{corollary:abelian-schemes} 
in \autoref{subsubsection:abelian-scheme-proof}.
In \autoref{proposition:hilbert-modular},
we prove a variant of the above corollary
with a stronger bound on the genus
when the abelian scheme has real multiplication, corresponding to a map from 
$C\setminus\{x_1, \cdots, x_n\}$ to a Hilbert modular stack.

\begin{remark}
\source{geometric-local-systems-general-curves-isomonodromy}
\notready
	\label{remark:remove-analytic}
\source{geometric-local-systems-general-curves-isomonodromy}
	Najmuddin Fakhruddin has pointed out to us that one may replace ``analytically very general" with "very general" 
	in the usual algebraic sense in \autoref{corollary:abelian-schemes}. Indeed, if a very general curve carried a non-isotrivial Abelian scheme of relative dimension less than $\sqrt{g+1}$, the same would be true for an analytically very general curve, by spreading out. 
	We do not know how to analogously strengthen \autoref{theorem:very-general-VHS}---see  \autoref{question:non-abelian-Hodge}.
\end{remark}

\begin{remark}
\source{geometric-local-systems-general-curves-isomonodromy}
\notready
    It is a well-known conjecture that integral local systems underlying a polarizable variation of Hodge structure are of geometric origin---see e.g.~\cite[Conjecture 12.4]{simpson62hodge} for a precise statement. \autoref{theorem:very-general-VHS} verifies this conjecture for local systems of rank less than $2\sqrt{g+1}$ on a analytically very general $n$-pointed hyperbolic curve of genus $g$, as local systems with finite monodromy arise from geometry.
\end{remark}	
	
We are grateful to H\'el\`ene Esnault for pointing out the following consequence of \autoref{corollary:geometric-local-systems} to us. We let $$\mathscr{M}_{B,r}(C\setminus \{x_1, \cdots, x_n\}):=\on{Hom}(\pi_1(C\setminus\{x_1, \cdots, x_n\}), \on{GL}_r(\mathbb{C}))\sslash \on{GL}_r(\mathbb{C})$$ be the \emph{character variety} parametrizing conjugacy classes of semisimple representations of $\pi_1(C\setminus \{x_1, \cdots, x_n\})$ into $\on{GL}_r(\mathbb{C})$. See e.g.~\cite{sikora2012character} for a useful primer on character varieties.
\begin{corollary}
\source{geometric-local-systems-general-curves-isomonodromy}
\notready
    \label{corollary:non-density-of-geometric-local-systems}
\source{geometric-local-systems-general-curves-isomonodromy}
    Let $(C, x_1, \cdots, x_n)$ be an analytically very general hyperbolic $n$-pointed curve of genus $g$. Then if $1<r<2\sqrt{g+1}$, the local systems of geometric origin are not Zariski-dense in the character variety $\mathscr{M}_{B,r}(C\setminus \{x_1, \cdots, x_n\})$.
\end{corollary}

\begin{remark}
\source{geometric-local-systems-general-curves-isomonodromy}
\notready
     \autoref{corollary:non-density-of-geometric-local-systems} contradicts conjectures of Esnault-Kerz \cite[Conjecture 1.1]{esnault2021local} and Budur-Wang \cite[Conjecture 10.3.1]{budur2020absolute}, which imply the density of geometric local systems in the character variety of any smooth complex variety.
\end{remark}



We will prove \autoref{corollary:non-density-of-geometric-local-systems} in 
\autoref{subsubsection:non-density-proof}.

In what follows, we say a flat vector bundle has \emph{unitary monodromy} if the associated monodromy representation
$\rho:\pi_1(C)\to \on{GL}_n(\mathbb{C})$
has image with compact closure.
We will deduce the above results from
\autoref{theorem:isomonodromic-deformation-CVHS} below,
using that a discrete subset of the image of a unitary $\rho$ is finite.

\begin{theorem}
\source{geometric-local-systems-general-curves-isomonodromy}
\notready\label{theorem:isomonodromic-deformation-CVHS}
\source{geometric-local-systems-general-curves-isomonodromy}
Let $(C, x_1, \cdots, x_n)$ be an $n$-pointed hyperbolic curve of genus $g$.
Let $({E}, \nabla)$ be a flat vector bundle on
$C$ with $\on{rk}{E}<2\sqrt{g+1}$ and with regular singularities at the $x_i$. If an isomonodromic
deformation of ${(E, \nabla)}$ to an analytically general nearby $n$-pointed curve underlies a polarizable complex variation of Hodge structure, then $({E},\nabla)$ has unitary monodromy.
\end{theorem}

\subsection{Main results on isomonodromic deformations}
\label{subsection:isomonodromy-intro}
\source{geometric-local-systems-general-curves-isomonodromy}
As remarked in \autoref{subsection:overview}, the Hodge-theoretic results of
\autoref{subsection:hodge-intro} arise from an analysis of the Harder-Narasimhan filtrations of isomonodromic deformations of flat vector bundles on curves. Our first such result is a counterexample to 
\cite[Theorem 1.3]{BHH:logarithmic},
\cite[Theorem 1.3]{BHH:irregular}, and
\cite[Theorem 1.2]{BHH:parabolic}, which demonstrates that the situation is somewhat more complicated than was previously believed --- there exist irreducible flat vector bundles whose isomonodromic deformations are never semistable. 

Specifically, \cite{BHH:logarithmic} ask the following
question. 
\begin{question}[\protect{\cite[p. 123]{BHH:logarithmic}}]
\source{geometric-local-systems-general-curves-isomonodromy}
\notready
	\label{question:bhh}
\source{geometric-local-systems-general-curves-isomonodromy}
	Let $X$ be a smooth proper curve, and $D\subset X$ a reduced effective divisor. Given a flat vector bundle $(E, \nabla)$ on $X$, with regular singularities along $D$, let $(E', \nabla')$ be the isomonodromic deformation of $(E,\nabla)$ to an analytically general nearby curve $(X', D')$. Is $E'$ semistable?
\end{question}

The main claim of \cite{BHH:logarithmic} is that \autoref{question:bhh} has a positive answer if $(E,\nabla)$ has irreducible monodromy and the genus of $X$ is at least $2$. However, the following results answer \autoref{question:bhh} in the negative, even in this case. See \autoref{remark:bhh-error} for a discussion of the errors in previous claims that \autoref{question:bhh} had a positive answer.

We use 
$\mathscr{M}_{g,n}$ to denote the analytic moduli stack of smooth proper curves with
geometrically connected fibers and $n$ distinct marked points.
\begin{theorem}
\source{geometric-local-systems-general-curves-isomonodromy}
\notready
	\label{theorem:counterexample}
\source{geometric-local-systems-general-curves-isomonodromy}
Let $g\geq 2$ be an integer. There exists a vector bundle with flat connection
$(\mathscr{F}, \nabla)$ on $\mathscr{M}_{g,1}$ such that for each fiber $C$ of
the forgetful morphism $\mathscr{M}_{g,1} \to \mathscr{M}_g$, the restriction of $(\mathscr{F}, \nabla)$ to $C$
\begin{enumerate}
    \item has semisimple monodromy and
    \item is not semistable.
\end{enumerate}
\end{theorem}
%We prove \autoref{theorem:counterexample} in
%\autoref{subsubsection:proof-counterexample}.
%

We also have the following variant, where the vector bundle has irreducible
monodromy, instead of just semisimple monodromy.
\begin{corollary}
\source{geometric-local-systems-general-curves-isomonodromy}
\notready\label{corollary:counterexample}
\source{geometric-local-systems-general-curves-isomonodromy}
Let $C$ be a smooth projective curve of genus at least $2$. There exists an irreducible flat
vector bundle $(E, \nabla)$ on $C$, whose isomonodromic deformations to
a nearby curve are never semistable.
\end{corollary}
\begin{proof}
	The restriction $(\mathscr{F}, \nabla)|_C$ from \autoref{theorem:counterexample} provides a semisimple flat vector bundle, each of whose flat summands has degree zero; by \autoref{theorem:counterexample}(2), its
isomonodromic deformation to a nearby curve is never semistable. Hence one of
the irreducible summands of
$(\mathscr{F}, \nabla)|_C$
satisfies the statement of the corollary.
\end{proof}

In a positive direction, we have have the following result, showing 
that the 
isomonodromic deformation of any semisimple flat vector bundle to an
analytically general nearby curve is close to being semistable, and moreover it is
semistable if the rank is small.

\begin{theorem}
\source{geometric-local-systems-general-curves-isomonodromy}
\notready
	\label{theorem:hn-constraints}
\source{geometric-local-systems-general-curves-isomonodromy}
	Let $(C,D)$ be hyperbolic of genus $g$ and let
	$({E}, \nabla)$ be a flat
	vector bundle on $C$ with regular singularities along $D$,
	and irreducible
monodromy.
Suppose
 $(E',\nabla')$ 
is an isomonodromic deformation
of $({E}, \nabla)$ to an analytically general nearby curve,
with Harder-Narasimhan filtration $0 = (F')^0 \subset (F')^1 \subset \cdots \subset
(F')^m =
E'$. For $1 \leq i \leq m$, let $\mu_i$ denote the slope of
$\on{gr}^{i}_{HN}E' := (F')^i/(F')^{i-1}$.
Then the following two properties hold.
\begin{enumerate}
	\item If $E'$ is not semistable, then for every $0 < i < m$, there
	exists $j < i < k$ with $$\rk \on{gr}^{j+1}_{HN}E'\cdot \rk
	\on{gr}^k_{HN}E'\geq g+1.$$
\item We have $0<\mu_i-\mu_{i+1}\leq 1$ for all $i<m$.
\end{enumerate}
\end{theorem}

In other words, the consecutive associated graded pieces of the generic
Harder-Narasimhan filtration have slope differing by at most one, and, if there
are multiple pieces of the generic Harder-Narasimhan filtration, many of them must have large rank relative to $g$.  This result is due to Bolibruch \cite[Proposition 5.6]{bolibrukh1990riemann} in the case that $g=0$ and $\on{rk}(E)=2$.
%We prove \autoref{theorem:hn-constraints} in \autoref{subsection:proof-hn}.

\autoref{theorem:hn-constraints} is a special case of
\autoref{theorem:hn-constraints-parabolic} below, where we allow certain
parabolic structures on the vector bundle $E$. These more general results are required for our Hodge-theoretic applications.
\begin{remark}
\source{geometric-local-systems-general-curves-isomonodromy}
\notready
	\label{remark:}
\source{geometric-local-systems-general-curves-isomonodromy}
	\autoref{theorem:hn-constraints} also holds without the hyperbolicity
	assumption, as we will explain. Nevertheless, it is convenient to make the assumption 
	so that curves have no infinitesimal automorphisms. In this case
	isomonodromic deformations are somewhat better behaved, see 
	\cite[p. 518]{Heu:universal-isomonodromic}.

	We now explain the proof of \autoref{theorem:hn-constraints} in the case
	$(C, D)$ is not hyperbolic.
	Suppose $(C, D)$ is not hyperbolic, so either $g = 1, n = 0$ or $g = 0, n \leq 2$. In
this case the fundamental group $\pi_1(C - \{x_1, \ldots, x_n\})$ is abelian.
This implies any irreducible representation of $\pi_1(C - \{x_1, \ldots, x_n\})$
is $1$-dimensional, so the corresponding flat vector bundle is a line bundle. In
this case, $E$ and $E'$ are semistable, so
\autoref{theorem:hn-constraints} still holds.
\end{remark}


As a corollary, we are able to salvage the main theorem of \cite{BHH:logarithmic} for flat vector bundles whose rank is small relative to $g$, using the AM-GM inequality.
The following corollary can be deduced directly from
\autoref{theorem:hn-constraints-parabolic} and AM-GM. 
It is also a special
case of \autoref{cor:stable-parabolic}.
\begin{corollary}
\source{geometric-local-systems-general-curves-isomonodromy}
\notready\label{cor:stable}
\source{geometric-local-systems-general-curves-isomonodromy}
	Let $(C, D)$ be a hyperbolic curve of genus $g$.
Let $({E}, \nabla)$ be an irreducible flat
vector bundle on $C$ with regular singularities along $D$, and suppose that
$\on{rk}(E)<2\sqrt{g+1}$. 
Then an isomonodromic deformation of 
$E$
to an analytically general nearby curve is  semistable. 
\end{corollary}
%We prove \autoref{cor:stable} in \autoref{subsection:main-stable}.
 As remarked above, \cite{BHH:logarithmic} claims an analogous theorem with no bound on the rank of $E$; our \autoref{corollary:counterexample} implies such a bound is necessary. Our methods are heavily inspired by those of \cite{BHH:logarithmic}, but our
technique requires some new input from Clifford's theorem for vector bundles.

These results appear to be new even for vector bundles with finite monodromy; we give some example applications in this case.
\begin{example}
\source{geometric-local-systems-general-curves-isomonodromy}
\notready
	\label{example:splitting-type}
\source{geometric-local-systems-general-curves-isomonodromy}
	In this example, we describe what 
	\autoref{theorem:hn-constraints}(2) tells us about splitting types
	of certain vector bundles on $\mathbb P^1$.
	Suppose we are given a finite group $G$ and
	a finite $G$-cover $f: X \to \mathbb P^1$.
	Consider the flat vector bundle $E := f_* \mathscr O_X$ on $\mathbb P^1$, with the connection $\nabla$ induced by the exterior derivative $d: \mathscr{O}_X\to \Omega^1_X$,
	which has regular singularities along the branch locus of $f$.
	Let $F$ be a summand of $E$ with irreducible monodromy
	and let $F'$ be an isomonodromic deformation of $F$ to a very general
	nearby pointed genus $0$ curve.
	Since every vector bundle on $\mathbb P^1$ is a sum of line bundles,
	we can write $F' = \mathscr O_{\mathbb P^1}(a_1)^{b_1} \oplus \cdots
	\oplus \mathscr O_{\mathbb P^1}(a_m)^{b_m}$,
	with $a_1 < a_2 < \cdots < a_m$ and $b_i > 0$.
	Then \autoref{theorem:hn-constraints}(1) tells us nothing, but
	\autoref{theorem:hn-constraints}(2) tells us that the $a_i$ are
	consecutive, i.e., $a_{i + 1} = a_i + 1$ for $1 \leq i \leq m-1$.

	Such $F'$ appear as a summand in $f'_* \mathscr O_{X'}$, where $f': X'
	\to \mathbb P^1$ is a general $G$-cover of $\mathbb P^1$.
\end{example}

\begin{example}
\source{geometric-local-systems-general-curves-isomonodromy}
\notready
	\label{example:tschirnhausen}
\source{geometric-local-systems-general-curves-isomonodromy}
	We now give a sample application of \autoref{cor:stable} to
	semistability of Tschirnhausen bundles of finite covers.
	Consider a family $\mathscr X \xrightarrow{\alpha} \mathscr{Z}
	\xrightarrow{\beta} \mathscr Y
	\xrightarrow{\gamma} B$
	where $\mathscr X \to B, \mathscr Y \to B$ and $\mathscr Z \to B$ are smooth proper curves
	with geometrically connected fibers,
	$\beta$ is finite locally free of degree $d$, $\beta\circ \alpha$ is an $S_d$ cover
	which is the Galois closure of $\beta$. 
	Suppose further $\beta \circ \alpha$ is branched
	over a divisor $\mathscr D \subset \mathscr Y$ consisting of $n$ disjoint sections over $B$ so that
	$(\mathscr Y, \mathscr D)$ is a relative hyperbolic curve of genus $g$ over $B$.
	Assume the map $B \to \mathscr M_{g,n}$ induced by $\gamma$
	is dominant.

	We obtain a flat vector bundle $E := (\beta \circ \alpha)_* \mathscr
	O_{\mathscr X}$ on $\mathscr Y$ with regular singularities along
	$\mathscr D$.
	One can decompose $E = \bigoplus\limits_{S_d \text{-irreps }
	\rho} E_\rho^{\dim \rho}$ as a sum of flat bundles with irreducible
	monodromy. Let $F$ denote one such
	summand corresponding to the standard representation of dimension $d-1$. 
	The flat vector bundle $\beta_* \mathscr O_{\mathscr Z}$ decomposes as
	$\mathscr O_{\mathscr Y} \oplus F$.
	The dual $T$ of $F$ is known
	as the Tschirnhausen bundle.
	By \autoref{cor:stable},
	the restriction of $T$ to a general fiber of $\mathscr Y \to B$
	is semistable whenever $d -1 < 2 \sqrt{g+1}$.

	Previous results on stability of $T$ were established in 
	\cite[Theorem 1.5]{deopurkarP:vector-bundles-and-finite-covers}.
	If $h$ denotes the genus of $\mathscr Z \to B$, they proved the
	restriction of $T$ to a general fiber of $\mathscr Y \to B$ is
	semistable whenever $h \geq dg + d(d-1)^2 g$
	\cite[Remark 3.16]{deopurkarP:vector-bundles-and-finite-covers}. After this paper appeared on the arXiv, it was shown by Coskun, Larson, and Vogt \cite{coskun2022stability} that the restriction of $T$ to a general fiber of $\mathscr Y \to B$
	is semistable when $g\geq 1$, and stable if $g\geq 2$.
\end{example}



\subsection{Motivation}
 Our main motivation comes from the following question. Let $f: X\to Y$ be a map
 of algebraic varieties. What are the restrictions on the topology of $f$? Our
 \autoref{corollary:geometric-local-systems} places a very strong restriction on
 the topology of morphisms to an analytically very general curve $C$ of genus
 $g$. For example, it implies that if $f: X\to C$ is a proper morphism with smooth generic fiber and bad reduction at $n$ analytically very general points of $C$, then any non-isotrivial monodromy representation occurring in the cohomology of $X/C$ has dimension at least $2\sqrt{g+1}$.
 
We became interested in this question and its connection to isomonodromy while
trying to understand \cite{BHH:logarithmic}. In that paper Biswas, Heu, and
Hurtubise raise \autoref{question:bhh}, asking whether it is possible to isomonodromically deform irreducible flat vector bundles to achieve semistability, by analogy to Hilbert's 21st problem (also known as the Riemann-Hilbert problem). 

Hilbert's 21st problem, as answered by  Bolibruch \cite{bolibruch1995riemann} (correcting earlier work of Plemelj), poses the question of whether every monodromy representation can be realized by a Fuchsian system. Esnault and Viehweg generalize this question to higher genus in \cite{esnault1999semistable}: they ask when an irreducible representation can be realized as the monodromy of a flat vector bundle $(E, \nabla)$ with regular singularities at infinity, with $E$ semistable.  
 
 In Esnault-Viehweg's formulation, the complex structure on the underlying curve
 is fixed, and the residues of the differential equation at regular singular
 points are modified to achieve semistability. Flipping this around,
 Biswas-Heu-Hurtubise's analogue asks if semistability can be achieved by
 modifying the complex structure and fixing the residues. They claim that this
 is always possible in the logarithmic, parabolic, and irregular settings, in
 \cite{BHH:logarithmic, BHH:parabolic, BHH:irregular}. After discovering the
 Hodge-theoretic counterexample to these claims in
 \autoref{corollary:counterexample}, we proved \autoref{theorem:hn-constraints}
 as an attempt  (1) to understand to what extent Biswas, Heu, and Hurtubise's
 \autoref{question:bhh} has a positive answer, and (2) to apply the cases when there is a positive answer to the analysis of variations of Hodge structure on curves.

\subsection{Idea of proof}
\label{subsection:idea-of-proof}
\source{geometric-local-systems-general-curves-isomonodromy}

To prove \autoref{theorem:very-general-VHS}, we first reduce to proving
\autoref{theorem:isomonodromic-deformation-CVHS}, using that discrete compact
spaces are finite.
We then prove \autoref{theorem:isomonodromic-deformation-CVHS} by showing that
any flat vector bundle satisfying the hypotheses of the theorem is forced to be (parabolically) semistable on an analytically general curve, whence the Hodge
filtration consists of a single piece by \autoref{corollary:unstable}.
The polarization then gives a definite Hermitian form preserved by the monodromy, and hence
the monodromy is unitary.
The key issue, which follows from \autoref{theorem:hn-constraints-parabolic},
is therefore to show that low rank flat vector bundles %underlying polarizable complex variations of Hodge structure 
are parabolically semistable
on an analytically general curve.

To prove \autoref{theorem:hn-constraints-parabolic}, we assume we have a flat
vector bundle $(E, \nabla)$ on our hyperbolic curve $(C, D)$, and consider an
isomonodromic deformation to a nearby curve.
To this end, we use the deformation theory of this flat vector bundle with its Harder-Narasimhan filtration,
which is governed by a variant of the Atiyah bundle.
We show that if the Harder-Narasimhan filtration does not satisfy the conclusion of \autoref{theorem:hn-constraints-parabolic}, then there is a direction along which we can deform the curve so that
the filtration is destroyed.
Indeed, if the filtration persisted, deformation theory provides us with a map
from $T_C(-D)$ to a certain
parabolically semistable subquotient of $\mathrm{End}(E)$ which vanishes on $H^1$.
Taking the Serre duals gives a semistable coparabolic vector bundle of low rank
and large coparabolic slope
which is not generically globally generated.
In the end, we rule this out by a variant of Clifford's theorem for vector
bundles.


\subsection{Organization of the paper}
\label{subsection:organization}
\source{geometric-local-systems-general-curves-isomonodromy}
In \autoref{section:parabolic}, we review background on parabolic bundles.
In \autoref{section:deformation-theory}, we give background on Atiyah bundles,
parabolic Atiyah bundles, and isomonodromic deformations. 
In \autoref{section:hodge-theoretic-preliminaries} we give background on complex variations of Hodge structures and their associated parabolic Higgs bundles. Experts can likely skip these three sections. 
In \autoref{section:counterexample}, we prove \autoref{theorem:counterexample} and \autoref{corollary:counterexample}, providing counterexamples to earlier published claims about semistability of isomonodromic deformations. 
In \autoref{section:hn-filtration} we prove the main results on isomonodromic
deformations, \autoref{theorem:hn-constraints} and \autoref{cor:stable} (and their generalizations \autoref{theorem:hn-constraints-parabolic} and \autoref{cor:stable-parabolic}). This is the technical heart of the paper. 
In \autoref{section:hodge-theoretic-results}, we prove the main consequences for
variations of Hodge structure, \autoref{theorem:very-general-VHS}, \autoref{corollary:geometric-local-systems}, \autoref{corollary:non-density-of-geometric-local-systems}, and \autoref{theorem:isomonodromic-deformation-CVHS}. 
Finally, \autoref{section:questions} lists some questions motivated by our results.

For readers unfamiliar with the theory of parabolic bundles, we suggest first considering the case of \autoref{theorem:very-general-VHS} where $\mathbb{V}$ is assumed to have unipotent monodromy at infinity. In this case one may replace parabolic stability with the usual notion of stability for vector bundles,  simplifying the proof;  in particular, one may use \autoref{cor:stable} in place of \autoref{cor:stable-parabolic}.
\subsection{Acknowledgments}
This material is based upon work supported by the Swedish Research Council under
grant no. 2016-06596 while the authors were in residence at Institut
Mittag-Leffler in Djursholm, Sweden during the fall of 2021. 
Landesman was supported by the National Science
Foundation under Award No. DMS-2102955;
Litt was supported by NSF grant DMS-2001196 and an NSERC grant, “Anabelian methods in arithmetic and algebraic geometry.” 
We would like to thank multiple anonymous referees for helpful comments.
We are grateful for useful conversations with 
Donu Arapura,
Indranil Biswas, 
Marco Boggi, 
Juliette Bruce, 
Anand Deopurkar,
H\'el\`ene Esnault, 
Najmuddin Fakhruddin,
Joe Harris,
Viktoria Heu,
Jacques Hurtubise, 
David Jensen, 
Jean Kieffer,
Matt Kerr,
Eric Larson, 
Haohao Liu,
Anand Patel,
Alexander Petrov,
Andrew Putman,
Will Sawin,
Ravi Vakil,
and Isabel Vogt.
